\documentclass[instructions.tex]{subfiles}
\begin{document}
 
\chapter{Des bons confesseurs.}
 
\primo Si tous les confesseurs, dit saint Charles, faisaient leur devoir dans l’esprit de Dieu, et selon les règles de l’Église, que de crimes n’empêcheraient-ils pas ! et combien de gens, qui courent à leur perte, seraient sauvés !
 
Le bon médecin n’est pas celui qui permet à son malade tout ce qu’il veut ; mais celui qui s’applique à connoître le mal, qui donne des remèdes efficaces, qui guérit, et qui prévient la rechute. De même, le bon confesseur est celui qui tâche de connoître les habitudes et les attaches d’un pénitent, qui lui donne des avis salutaires, qui le convertit et le préserve de la rechute. Un médecin doux est souvent un mauvais médecin, qui laisse de la corruption dans les plaies, et du venin dans le corps. De même, les confesseurs lâches et complaisans laissent souvent les âmes dans le péché et dans la mort.
 
Il faut qu’un confesseur soit savant, et plus savant qu’on ne pense ; mais il faut aussi qu’il ait du zèle et de la fermeté. De quoi me sert-il d’avoir un médecin savant, s’il me laisse mourir. De quoi me sert-il que mon confesseur ait beaucoup de science, s’il me laisse damner ? Il y a bien plus de confesseurs trop humains et trop complaisans, que de confesseurs ignorans ; les uns ne sont pas moins pernicieux aux âmes que les autres.
 
\secundo Les confesseurs ont tous le pouvoir de retenir et de remettre les péchés, de lier et de délier les consciences ; ils doivent donc se servir avec prudence de ce double pouvoir. Ceux qui ne se servent presque jamais du pouvoir de retenir les péchés, se lient souvent eux-mêmes, en pensant délier les autres.
 
Certains remèdes guérissent les uns et tuent les autres, soulagent aujourd’hui et font du mal en un autre temps : de même l’absolution donne la vie à un pénitent bien disposé, et laisse dans la mort celui qui ne l’est pas. Les absolutions sont comme des lettres de grâce ; mais elles ne sont pas toutes ratifiées dans le ciel. Urie croyait porter à son général une lettre favorable, et il portait, sans le savoir, l’arrêt de sa mort. Il y a des criminels qui, croyant avoir une lettre de grâce, trouvent leur condamnation en la présentant au juge. Oh ! combien d’absolutions feront un jour la confusion des pénitens et des confesseurs !
 
Si le confesseur doit éviter le scrupule, et traiter ses pénitens avec beaucoup de charité et de douceur, il doit aussi leur parler avec autorité, parce qu’il tient la place de Jésus-Christ, et ne pas trahir les intérêts de son maître, en donnant l’absolution à des incorrigibles qui abusent des sacremens. S’il fait tort aux fidèles, en leur refusant les sacremens sans raison légitime, il ne leur fait pas moins de tort s’il les leur accorde sans dispositions suffisantes : en rebutant ses pénitens par une conduite dure et austère, il les désespère et leur ferme le ciel ; mais si, par respect humain, par une fausse compassion et par lâcheté, il les laisse vivre à leur liberté, il leur ouvre l’enfer.
 
Un confesseur qui vous absout, lorsque vous êtes dans l’occasion prochaine du péché mortel et dans l’habitude du vice, qui vous laisse dans un état de damnation, faute de vous prescrire les moyens de vous en retirer, est, selon la doctrine des saints Pères, un guide aveugle, un meurtrier de votre âme, un profanateur du sang de Jésus-Christ, et un dissipateur des saints mystères\footnote{Dissipatores, non dispensatores mysteriorum Christi. SS. PP.}.
 
\section{Résolutions}
 
\primo Quand je serai dans le cas de choisir un confesseur, je consulterai Dieu par la prière. 

\secundo J’en chercherai alors un qui réunisse à la science et à l’expérience dans la conduite des âmes, l’exacte observation des règles que lui prescrit son redoutable ministère. 

\tertio Et quand j’aurai eu le bonheur de réussir dans ce choix, je donnerai à ce confesseur une confiance sans bornes pour tout ce qui regarde mon salut. 

\prayer{Mon Dieu, aidez-moi à remplir tous mes devoirs envers mon guide spirituel.}
 
\end{document}
